\chapter{소자 수준 회로 모델(Simulink) 현황}

\section{테스트벤치 시리즈}

소자 수준 회로 모델은 검증 단계마다 새 테스트벤치를 만들어 남겼다. 모든 모델은 빌드 스크립트로 처음부터 생성된다. 각 테스트벤치는 같은 조건의 검증 모델 결과를 기준값으로 두어 시험 전에 정한 기준으로 판정하는 것이 원칙이다(표 \ref{t:tb}; 예외는 18.2절). 중간보고서(TB01–TB20) 이후 TB21–TB33을 더했고, 불통과 결과도 지우지 않았다.

그림 \ref{f:sl19}은 이 가운데 Stage A·B·C와 제어기를 한 회로에 넣은 TB19의 최상위 화면이다. 이 회로의 Stage A는 이전 150 kHz 설계이며, 최종 A25R Stage A 회로와 Stage B를 함께 넣은 테스트벤치는 없다.

\begin{figure}[!htb]
\centering
\includegraphics[width=16.0cm]{\repo/docs/fig/report/fig17_1_tb19_top.png}\par\smallskip
{\footnotesize\raggedright
\textbf{1} Stage A 모듈 1(이전 150 kHz 설계) \quad \textbf{2} Stage B: R2-W460 창 조정 모듈(바이패스 포함) \quad \textbf{3} 400 V 버스·부하(정전력 부하, DAB 손실 등가) \quad \textbf{4} 조정형 Stage C \quad \textbf{5} 48 V 출력단·부하 \quad \textbf{6} 게이트·제어·보호 블록(Gate\_StageA, Controller\_StageB, Gate\_StageB, Controller\_StageC, Gate\_StageC)\par}
\caption{TB19 최상위 회로(MATLAB/Simulink): R2-W460 통합 강등과 Stage A 경부하 위상 셰딩}\label{f:sl19}
\srcnote{주: TB19 Simulink 최상위 회로 화면에서 하단 설명을 제외하고 영역 번호를 덧붙였다. 인쇄 크기에서 화면 속 블록 이름은 읽기 어려우므로 영역은 위 범례로 구분한다. Stage A는 \textbf{이전 150 kHz 설계 1모듈}(최종 A25R 아님)이고, 48 V 쪽 배터리 가지는 1 MΩ으로 끊어 배터리가 없다. 보호 순서는 Stage C 일시 정지(9.6절) 도입 전이며, 판정한 것은 H-브리지 개방 강등 2사례(통과 2/2, 표 \ref{t:tb})다. 하드웨어가 아닌 Simscape 소자 수준 회로다.}
\end{figure}

{\footnotesize
\begin{longtable}{L{\dimexpr 0.10\linewidth-2\tabcolsep\relax}L{\dimexpr 0.40\linewidth-2\tabcolsep\relax}L{\dimexpr 0.11\linewidth-2\tabcolsep\relax}L{\dimexpr 0.39\linewidth-2\tabcolsep\relax}}
\caption{테스트벤치 시리즈와 판정}\label{t:tb}\\
\toprule
\tabhead \textbf{TB} & \textbf{대상} & \textbf{기준} & \textbf{판정} \\
\midrule
\endfirsthead
\toprule
\tabhead \textbf{TB} & \textbf{대상} & \textbf{기준} & \textbf{판정} \\
\midrule
\endhead
\bottomrule
\endlastfoot
TB01 & Stage A 1위상 RSC, 고정 타이밍 & E1 & 통과(전 지표 0.2 \% 이내) \\
TB02 & Stage B H-브리지(개루프), 데드타임 전압오차 & 교정표 & 통과(최대 차 0.135 V) \\
TB03 & Stage B DAB(개루프), 위상 계단·직류 편류 & E12 & 통과(0.39 A 이내) \\
TB04 & Stage B 폐루프 & E12 & 일치 통과(리플은 TB04B에서 결론) \\
TB04B & TB04 + 후보 소자 출력용량, 게이트 격자 & E12 & 통과(0.625 ns, 리플 0.069·0.031 V) \\
TB05 & Stage C CLLC 3 kW(고정이득) & E10 & 통과(이득 0.9965–1.0001) \\
TB06 & 통합(주행): Stage A 모듈 1 + B + C & E12, E1, TB05 & 통과 3/3 \\
TB07 & 통합(충전): Stage A 8위상 + 충전기 + B + C & E16, E1 & 회로 한계 7/7, 일치는 조건부(충전전류 램프 ≥ 100 µs) \\
TB08 & 창 조정 모듈 R2-W460 + 바이패스·모드 시퀀서 & E12, E19 & 통과 5/5 \\
TB09 & R2 보호 상태기계, 제동 저항, 소프트 투입, 감쇠기 & E17 & 최종 실행 4/4(첫 실행 2사례 불통과 → 상태기계 보완) \\
TB10 & DAB 고장 모드 변조(EPS 게이트, 정합 유지, 감발) & E20 & 통과 3/3(회로 기준 2배), 1.5배 목표는 붕괴 구간 초과 \\
TB11 & TB10 + 데드타임 반영 첨두 표 & E21 & 불통과(DAB 첨두 88.5·68.9 A) — 지령표 칸 전환 과도 \\
TB12 & TB11 + EPS 지령 전이의 직류 편류 방지·변화율 상한 & E21 & 1/2 통과(1.5배 기준), DAB 첨두 70.0·68.5 A — 전이 끝 1.3 A 초과 \\
TB13 & TB12 + 지령 전이 여유(설계 상한 1.387배) & E21 & 통과 2/2(1.5배 기준), DAB 첨두 66.4·64.4 A \\
TB14 & 조정형 Stage C 5 kW(동기 벅 + CLLC 직류변압기 4.8:1) & E22 & 1차 0/5(영전압 스위칭) → 재설계 뒤 2차 6/7(270 V 역방향 경부하 불통과) \\
TB15 & TB14 + 경부하 두 브리지 구동 & E22 & 통과 5/5 \\
TB16 & R2 기본형 통합: Stage A 모듈 1 + 버스 감쇠기 + 조정형 Stage C & E23, E1 & 통과 3/3 \\
TB17 & R2-W460 통합 강등: TB13 + Stage A·조정형 Stage C 회로 & E23, E17 & 통과 2/2(강등 중 48 V ±0.03 \%) \\
TB18 & Stage A 경부하 위상 셰딩: 모듈 1 + R2 버스 + 충전기 전류원 & E24 & 통과 5/5 \\
TB19 & TB17 + 위상 셰딩 게이트 & E24, E17 & 통과 2/2 \\
TB20 & 충전 중 정지: Stage A 4모듈 + 충전기 + V\textsubscript{A} 클램프 + 버스 + 조정형 Stage C & E26 & 3/4(충전기 램프 사례 턴오프 1.01 A > 1.0 A) \\
TB21 & TB20 + 영교차 게이트 방향 반전 처리 & E27 & 통과 4/4(사례 D 턴오프 최대 0.83 A) \\
TB22 & Stage A 4모듈 관리기 + 위상 셰딩(150 kHz) & E28 & 3/4(사전 기동 없는 충전 불통과, 모델과 같음) \\
TB23 & TB13 + SALS 부하 블록·상태기계 연동(SALS 시간 0.15배) & E29 & 판정이 모델과 5/5 같음(통과 4, 회생 차단을 뺀 비교 사례 불통과) \\
TB24 & Stage A 4모듈 모듈 고장 처리 & E30 & 통과 4/4 \\
TB25 & Stage A 1위상 + 스위치 등가 출력용량(150·75·25 kHz) & E31 & 무부하 출력용량 손실 3/3, 부하 대조 3/6(T2 2/3·T3 1/3, 고정 타이밍 게이트·적분 간격), 전체 6/9 \\
TB26 & A25 통합(25 kHz, C\textsubscript{A} 400 µF) & E33 & 3/3 불통과(동특성 판정은 모델과 같음, 손실 대조 방식 불일치) \\
TB27 & A25R 통합(C\textsubscript{A} 2.4 mF) & E35 & 통과 4/4(FF27·MC27·MA27·SC27; SC27 한 사례의 손실 회로 대조 +3.0 \%) \\
TB28 & A25R 충전 중 정지·V\textsubscript{A} 클램프(1 Ω) & E26, E36 & 통과 4/4(경로 상실 A28·B28, 모듈 고장 C28, 충전 램프 D28) \\
TB29 & TB28 + 클램프 1.2 Ω & E37 & 통과 2/2 \\
TB30 & TB28 + 클램프 가지 2개 × 1.6 Ω, 저항 모서리 & E38 & 통과 2/2(버스 469.35 V) \\
TB31 & TB23 + C\textsubscript{A} 2.4 mF + Stage C 일시 정지 + 400 V 부하 정지 & E38 & 통과 2/2 \\
TB32 & TB30 + 가지 저항 섞임 & E39 & 통과 2/2(작은 저항 가지 393.0 A) \\
TB33 & 최종 클램프(가지 3개 × 2.4 Ω) + 최대 연속 켜짐 0.5 ms & E26, E40 & 경로 상실 2/2 통과, 오동작 사례 불통과(평균모델 대조 항목, 회로가 9 µs 만에 끔). 진단 2건(판정 아님) \\
\end{longtable}
}

\section{모델 간 차이와 해석상 유의점}

\begin{itemize}
\item \textbf{게이트 에지 양자화.} 게이트 에지가 적분 간격에 맞춰지며, 이것이 리밋 사이클 리플의 원인이었다(TB04·TB04B). 0.625 ns 이하에서 기준을 만족한다. TB06–TB20은 2.5 ns(TB06은 5 ns)로 판정했다.
\item \textbf{Stage A 등가와 설계.} TB08–TB13·TB23·TB31의 V\textsubscript{A}는 이상 전원(V\textsubscript{bat}/2 + 13.2 mΩ + 1.5 µH; TB31은 C\textsubscript{A} 2.4 mF 추가)이고, TB16–TB20에서 150 kHz Stage A 회로로 등가를 확인했다(공진 주기 약 30 µs보다 느린 외란). Stage A는 TB01–TB24가 150 kHz, TB26이 A25, TB27–TB33이 A25R이다.
\item \textbf{바이패스와 소프트 투입.} 루프의 기생 인덕턴스와 셀 출력용량은 넣지 않았고(투입 펄스는 상한 쪽 값), 소프트 투입은 가변 저항으로 모델링해 소자 온도는 모델 밖이다.
\item \textbf{해석기.} 후진 오일러는 공진망에 수치 감쇠를 만든다(TB05). 조정형 Stage C 레그 공진도 감쇠되지만 가변 간격 해석과의 대조(TB14 사례 A)에서 10 µs 평균 차는 없었다. DAB 손실은 회로 저항분만 남는다.
\item \textbf{V\textsubscript{A} 클램프.} 같은 게이트로 켜지는 가지들은 등가 가지로 묶었다(TB30 하나, TB32·TB33은 작은 저항 가지 하나와 나머지 등가). 클램프 스위치의 선정식 환산온도는 회로가 아니라 평균모델 사건 손실과 데이터시트로 계산했다(16장).
\end{itemize}
